\section{BRS 对称性}\label{sec:brs}

现在开始一般性讨论，它将最终导出体场关联函数
按规范耦合展开至所有阶均有限的证明。作为第一步，
本节回顾 $D+1$ 维中理论的 BRS 对称性。

小节 \ref{ssec:gaugefixing} 讨论的规范固定遵循标准程序，因此给出一个具有
BRS 对称性的理论。该对称性照常作用在边界场上，
但体场的变换需要扩散方程 \eqref{eq:3.2} 的解，
因而往往是非定域的。不过，Zinn--Justin 与 Zwanziger 在关于
Langevin 方程重正化的工作 \cite{ZinnZwanziger} 中证明，
通过引入附加的鬼场可以恢复变换的定域性。

\subsection{边界场的 BRS 变换}

未重正化场 $A_{\mu}$、$\cbar$、$c$ 的 BRS 变分定义为 \cite{BRSI,BRSII}
\begin{align}
   \dbrs A_{\mu}=D_{\mu}c,
   \label{eq:6.1}\\[2pt]
   \dbrs c=-c^2,
   \label{eq:6.2}\\[2pt]
   \dbrs\cbar=\lambda_0\partial_{\mu}A_{\mu}.
   \label{eq:6.3}
\end{align}
注意，鬼场的分量 $c^a$、$\cbar^a$ 以及算符 $\dbrs$
彼此反对易。例如式 \eqref{eq:6.2} 右端的乘积为
\begin{equation}
   c^2=c^ac^bT^aT^b=\frac{1}{2}c^ac^bf^{abc}T^c,
   \label{eq:6.4}
\end{equation}
而且算符 $\dbrs$ 的 Leibniz 法则必须计及其反对易性质。
于是容易证明：$D$ 维中理论的作用量 $S+\Sgf+\Sgh$
与其泛函积分中的测度都是 BRS 不变的。

\subsection{体鬼场}

前面提到的附加鬼场 $d(t,x)$ 与 $\dbar(t,x)$ 生活在 $D+1$ 维中，
但除此之外与 Faddeev--Popov 鬼场同属一类。它们的作用量为
\begin{equation}
   \Sflgh=-2\int_{0}^{\infty}\rmd t\int\rmd^Dx\,
   \tr\bigl\{
   \dbar(t,x)\bigl(\partial_td-\alpha_0 D_{\mu}\partial_{\mu}d\bigr)(t,x)
   \bigr\},
   \label{eq:6.5}
\end{equation}
且要求 $d$ 场满足边界条件
\begin{equation}
   \left.d\right|_{t=0}=c.
   \label{eq:6.6}
\end{equation}
$\dbar$ 场不施加边界条件；在许多方面它的角色类似于规范场情形下的
拉格朗日乘子场 $L_{\mu}$。

鬼场的传播子为
\begin{align}
   \langle\dtilde^a(t,p)\dbartilde^b(s,q)\rangle=
   (2\pi)^D\delta(p+q)\delta^{ab}\theta(t-s)\tilde{K}_{t-s}(p)
   +\rmO(g_0^2),
   \label{eq:6.7}\\[3pt]
   \tilde{K}_{t}(p)=\rme^{-\alpha_0tp^2},
   \label{eq:6.8}
\end{align}
可以按照小节 \ref{ssec:propagators4} 的步骤算出。
图形上该传播子用一条虚线（dashed line）表示：
\begin{equation}
  \raisebox{-0.55cm}{\includegraphics[width=1.9cm]{images/prop4}}%
  \;=\;\delta^{ab}\theta(t-s)\tilde{K}_{t}(p),
  \label{eq:6.9}
\end{equation}
箭头沿从 $\dbar$ 指向 $d$ 的方向画出。
还存在一个混合传播子
\begin{align}
   \langle\dtilde^a(t,p)\cbartilde^b(q)\rangle=
   (2\pi)^D\delta(p+q)\delta^{ab}g_0^2\Dtilde_t(p)
   +\rmO(g_0^4),
   \label{eq:6.10}\\[3pt]
   \Dtilde_t(p)={1\over p^2}\rme^{-\alpha_0tp^2},
   \label{eq:6.11}
\end{align}
它在 $t=0$ 处与 $c\cbar$ 传播子重合，因而用同一图形符号
（带方向的点线）表示。另一方面，$c\dbar$ 两点函数按所有阶都为零。

作用量
\begin{align}
   \left.\Sflgh\right|_{\rm interaction}=
   -\int_0^\infty\rmd t\int_{p,q,r}(2\pi)^D\delta(p+q+r)
   \notag\\[2pt]
   \times\vbulk{1,1}(p,q,r)^{abc}_{\mu}
   \tilde{B}^a_{\mu}(t,-p)\dbartilde(t,-q)^b\dtilde(t,-r)^c
   \label{eq:6.12}
\end{align}
还产生一个新的流顶点 $\vbulk{1,1}$。在费曼图中，
它像其他流顶点一样用空心圆圈表示。顶点的显式表达式见附录~\ref{app:B}。

正如小节 \ref{ssec:flowlines} 中讨论的流线圈的情形，
可以证明带 $d\dbar$ 线圈的图为零。注意，含混合传播子的鬼圈并不存在，
因为没有 $c\dbar$ 传播子，而且各顶点只把 $c$ 与 $\cbar$ 相连
或把 $d$ 与 $\dbar$ 相连。于是唯一非零的鬼圈
就是流时间为零处理论的通常鬼圈。

\subsection{体场的 BRS 变分}

BRS 对称性按如下方式作用于体场 \cite{ZinnZwanziger}：
\begin{align}
   \dbrs B_{\mu}=D_{\mu}d,
   \label{eq:6.13}\\[2pt]
   \dbrs L_{\mu}=[L_{\mu},d],
   \label{eq:6.14}\\[2pt]
   \dbrs d=-d^2,
   \label{eq:6.15}\\[2pt]
   \dbrs\dbar=D_{\mu}L_{\mu}-\{d,\dbar\}.
   \label{eq:6.16}
\end{align}
注意，在边界上 $\dbrs B_{\mu}=\dbrs A_{\mu}$ 且 $\dbrs d=\dbrs c$，
正如由边界条件 \eqref{eq:2.3} 与 \eqref{eq:6.6} 所应有的那样。

为了证明体作用量 $\Sfl+\Sflgh$ 的 BRS 变分为零，
引入场
\begin{align}
   E_{\mu}=\partial_tB_{\mu}-D_{\nu}G_{\nu\mu}
   -\alpha_0D_{\mu}\partial_{\nu}B_{\nu},
   \label{eq:6.17}\\[3pt]
   e=\partial_td-\alpha_0 D_{\mu}\partial_{\mu}d,
   \label{eq:6.18}
\end{align}
会有帮助。经过一些代数运算可以发现
\begin{align}
   \dbrs E_{\mu}=[E_{\mu},d]+D_{\mu}e,
   \label{eq:6.19}\\[3pt]
   \dbrs e=-\{e,d\},
   \label{eq:6.20}
\end{align}
于是体作用量的不变性便容易确立了。

\begin{figure}[t]
\centerline{\includegraphics[width=10.0cm]{images/diagram9}}
\caption{对式 \eqref{eq:6.23} 左端的关联函数有贡献的树图（图 1 与图 2）
以及对右端关联函数有贡献的树图（图 3、4、5）。
动量全部取为入射并满足 $p+q+r=0$。在图 $3$ 中，
动量 $p$ 流入一个代表场 $[B_{\mu},d]$ 插入的顶点；
而对图 $4$ 与图 $5$，带动量 $p$ 的外线要乘上 $ip_{\mu}$
（从而代表场 $\partial_{\mu}d$）。}
\label{fig:brsdiag}
\end{figure}

\subsection{BRS 恒等式的例子}

由于泛函积分测度也是不变的，故对基本场的任意乘积 $\mathcal{O}$
都有 $\langle\dbrs\mathcal{O}\rangle=0$。由此得到诸如
\begin{align}
   \lambda_0\langle
   \partial_{\mu}A_{\mu}^a(x)\partial_{\nu}A_{\nu}^b(y)
   \rangle=g_0^2\delta^{ab}\delta(x-y),
   \label{eq:6.21}\\[3pt]
   \lambda_0\langle B^a_{\mu}(t,x)\partial_{\nu}A^b_{\nu}(y)\rangle=
   -\langle(D_{\mu}d)^a(t,x)\cbar^b(y)\rangle,
   \label{eq:6.22}\\[3pt]
   \lambda_0\langle
   B^a_{\mu}(t,x)\partial_{\nu}A^b_{\nu}(y)\partial_{\rho}A^c_{\rho}(z)
   \rangle=
   -\langle(D_{\mu}d)^a(t,x)\cbar^b(y)\partial_{\rho}A^c_{\rho}(z)\rangle
   \label{eq:6.23}
\end{align}
之类的恒等式，其中第一个就是 $D$ 维理论中人们熟悉的规范 Ward 恒等式，
它确定规范场两点函数的纵向部分。

在微扰论的树级检验式 \eqref{eq:6.22} 与 \eqref{eq:6.23} 是有启发性的。
第一式把混合规范场传播子的纵向部分与混合鬼传播子联系起来，
这个恒等式从传播子的表达式看是显然的。
另一式涉及 $3$ 点顶点（见图~\ref{fig:brsdiag}）
以及各贡献之间的一次不平凡的相消（见附录~\ref{app:C}）。
